\documentclass[a4paper]{article}
\usepackage{ctex}
\ctexset{proofname=\heiti{证明}}
\usepackage{amsmath,amssymb,amsthm}
\usepackage{mathtools}
\usepackage{bm}
\usepackage{enumitem}
\usepackage{booktabs}
\usepackage{array}
\usepackage{graphicx}

\IfFileExists{iidef.sty}{
    \usepackage[thehwcnt = 8]{iidef}
    \thecourseinstitute{清华大学}
    \thecoursename{人工智能原理}
    \theterm{2026年春季学期}
    \hwname{作业}
    \slname{\heiti{解}}
}{
    \makeatletter
    \newcommand{\thecourseinstitute}[1]{\def\@courseinstitute{#1}}
    \newcommand{\thecoursename}[1]{\def\@coursename{#1}}
    \newcommand{\theterm}[1]{\def\@term{#1}}
    \newcommand{\hwname}[1]{\def\@hwname{#1}}
    \newcommand{\slname}[1]{\def\@slname{#1}}
    \thecourseinstitute{清华大学}
    \thecoursename{人工智能原理}
    \theterm{2026年春季学期}
    \hwname{作业}
    \slname{\heiti{解}}
    \newcommand{\courseheader}{
        \begin{center}
            {\Large \heiti \@courseinstitute \quad \@coursename}\\[0.5em]
            {\large \@term \quad \@hwname 8 \quad \@slname}
        \end{center}
        \vspace{1em}
    }
    \newcommand{\name}[1]{\noindent #1 \par\vspace{1em}}
    \newenvironment{solution}{\par\noindent\textbf{解：}\par}{\par}
    \makeatother
}

\begin{document}
\courseheader
\name{姓名：杜一郎\qquad 学号：2023011618}

\begin{enumerate}
\setlength{\itemsep}{3\parskip}

\item 两状态 MDP

状态集合 \(\mathcal S=\{s_0,s_1\}\)，动作集合
\(\mathcal A=\{a_0,a_1\}\)，折扣因子 \(\gamma=0.9\)。

\begin{solution}
\textbf{(1)}

策略为
\[
\pi(s_0)=a_1,\qquad \pi(s_1)=a_0,
\]
题目中还给出
\[
U^\pi(s_0)=5.
\]
状态 \(s_1\) 下执行 \(a_0\)：以 \(0.8\) 概率转移到 \(s_0\)，奖励为 \(2\)；以 \(0.2\) 概率转移到 \(s_1\)，奖励为 \(0\)。Bellman 方程为
\[
\begin{aligned}
U^\pi(s_1)
&=0.8\left(2+\gamma U^\pi(s_0)\right)
  +0.2\left(0+\gamma U^\pi(s_1)\right)\\
&=0.8(2+0.9\times 5)+0.2\times0.9\,U^\pi(s_1).
\end{aligned}
\]
整理得
\[
(1-0.18)U^\pi(s_1)=5.2,
\]
解得
\[
U^\pi(s_1)=\frac{5.2}{0.82}\approx 6.3415.
\]

\textbf{(2)}

看 \(s_0\) 改选 \(a_0\) 时的动作价值。执行 \(a_0\) 时，以 \(0.7\) 概率转移到 \(s_0\)，奖励为 \(1\)；以 \(0.3\) 概率转移到 \(s_1\)，奖励为 \(3\)。于是
\[
\begin{aligned}
Q^\pi(s_0,a_0)
&=0.7\left(1+\gamma U^\pi(s_0)\right)
 +0.3\left(3+\gamma U^\pi(s_1)\right)\\
&=0.7(1+0.9\times5)
 +0.3(3+0.9\times6.3415)\\
&\approx 6.4622.
\end{aligned}
\]
这里
\[
Q^\pi(s_0,a_0)\approx 6.4622 > U^\pi(s_0)=5,
\]
\(s_0\) 处把策略从 \(a_1\) 改为 \(a_0\) 会提高价值，即
\[
\pi_{\mathrm{new}}(s_0)=a_0.
\]
\end{solution}

\item \(3\times3\) Gridworld

状态编号如下：
\[
\begin{matrix}
s_1&s_2&s_3\\
s_4&s_5&s_6\\
s_7&s_8&s_9
\end{matrix}
\]
其中 \(s_9\) 为终止状态，\(V(s_9)=0\)。动作集合为
\[
\mathcal A=\{\mathrm{Up},\mathrm{Down},\mathrm{Left},\mathrm{Right}\}.
\]
每走一步奖励 \(R=-1\)，折扣因子 \(\gamma=1\)。若动作撞墙，则停留在原状态。初始价值为
\[
V_0(s)=0,\qquad s\in\{1,2,\ldots,9\},
\]
初始策略为均匀随机策略
\[
\pi_0(a|s)=0.25.
\]

\begin{solution}
\textbf{(1)}

第一次迭代用
\[
V_{k+1}(s)=\sum_a \pi_0(a|s)\left[-1+V_k(s')\right].
\]
由 \(V_0(s)=0\)，所有非终止状态都有
\[
V_1(s)=-1,\qquad s=1,2,\ldots,8.
\]
终止状态仍为
\[
V_1(s_9)=0.
\]

\textbf{(2)}

第二次迭代时，非终止后继状态的 \(V_1\) 为 \(-1\)，终止状态 \(s_9\) 的价值为 \(0\)。

对 \(s_5\)，四个方向分别到达 \(s_2,s_8,s_4,s_6\)，都不是终止状态，故
\[
V_2(s_5)=\frac14[(-1-1)+(-1-1)+(-1-1)+(-1-1)]=-2.
\]

对 \(s_6\)，向上到 \(s_3\)，向下到 \(s_9\)，向左到 \(s_5\)，向右撞墙仍在 \(s_6\)。代入得
\[
V_2(s_6)=\frac14[(-1-1)+(-1+0)+(-1-1)+(-1-1)]
=-\frac74=-1.75.
\]

对 \(s_8\)，向上到 \(s_5\)，向下撞墙仍在 \(s_8\)，向左到 \(s_7\)，向右到 \(s_9\)。代入得
\[
V_2(s_8)=\frac14[(-1-1)+(-1-1)+(-1-1)+(-1+0)]
=-\frac74=-1.75.
\]

\textbf{(3)}

用 \(V_2\) 做策略改进：
\[
\pi_1(s)=\arg\max_a\left[-1+V_2(s')\right].
\]
用到的几个值是
\[
\begin{gathered}
V_2(s_2)=V_2(s_3)=V_2(s_4)=V_2(s_5)=-2,\\
V_2(s_6)=V_2(s_8)=-1.75,\\
V_2(s_9)=0.
\end{gathered}
\]

在 \(s_5\) 处：
\[
\begin{array}{c|c}
\toprule
\text{动作} & -1+V_2(s')\\
\midrule
\mathrm{Up} & -3\\
\mathrm{Down} & -2.75\\
\mathrm{Left} & -3\\
\mathrm{Right} & -2.75\\
\bottomrule
\end{array}
\]
最优动作有两个：
\[
\pi_1(s_5)\in\{\mathrm{Down},\mathrm{Right}\}.
\]

在 \(s_6\) 处：
\[
\begin{array}{c|c}
\toprule
\text{动作} & -1+V_2(s')\\
\midrule
\mathrm{Up} & -3\\
\mathrm{Down} & -1\\
\mathrm{Left} & -3\\
\mathrm{Right} & -2.75\\
\bottomrule
\end{array}
\]
最优动作为
\[
\pi_1(s_6)=\mathrm{Down}.
\]
\end{solution}

\item Monte Carlo 价值估计

状态集合 \(\mathcal S=\{A,B\}\)，终止状态 \(T\)，折扣因子
\(\gamma=1\)。已采样得到三条轨迹：
\[
\begin{aligned}
&1:\quad A \xrightarrow{R=-1} B \xrightarrow{R=-1} A
\xrightarrow{R=-2} B \xrightarrow{R=+6} T,\\
&2:\quad B \xrightarrow{R=-2} B \xrightarrow{R=-1} A
\xrightarrow{R=+5} T,\\
&3:\quad A \xrightarrow{R=-1} A \xrightarrow{R=-3} B
\xrightarrow{R=+4} T.
\end{aligned}
\]

\begin{solution}
\textbf{(1) First-visit Monte Carlo}

先算每条轨迹中每个状态第一次出现时的回报。

第 1 条轨迹：
\[
G(A)= -1-1-2+6=2,\qquad
G(B)= -1-2+6=3.
\]
第 2 条轨迹：
\[
G(B)= -2-1+5=2,\qquad
G(A)=5.
\]
第 3 条轨迹：
\[
G(A)= -1-3+4=0,\qquad
G(B)=4.
\]
first-visit 估计为
\[
V(A)=\frac{2+5+0}{3}=\frac73\approx 2.3333,
\]
\[
V(B)=\frac{3+2+4}{3}=3.
\]

\textbf{(2) Every-visit Monte Carlo}

every-visit 要把同一条轨迹里同一状态的多次出现都计入平均。

状态 \(A\) 的所有回报为
\[
2,\ 4,\ 5,\ 0,\ 1.
\]
于是
\[
V(A)=\frac{2+4+5+0+1}{5}=\frac{12}{5}=2.4.
\]

状态 \(B\) 的所有回报为
\[
3,\ 6,\ 2,\ 4,\ 4.
\]
于是
\[
V(B)=\frac{3+6+2+4+4}{5}=\frac{19}{5}=3.8.
\]
\end{solution}

\end{enumerate}
\end{document}
